\documentclass[11pt]{article}
\usepackage[margin=1in]{geometry}
\usepackage{amsmath,amssymb,amsthm,mathtools}
\usepackage{hyperref}
\usepackage{booktabs}

\newtheorem{theorem}{Theorem}
\newtheorem{lemma}[theorem]{Lemma}
\newtheorem{proposition}[theorem]{Proposition}
\newtheorem{corollary}[theorem]{Corollary}
\newcommand{\E}{\mathbb E}
\newcommand{\Pp}{\mathbb P}
\newcommand{\Var}{\operatorname{Var}}
\newcommand{\N}{\mathcal N}

\title{Cover-Time Fluctuations Can Exceed the Maximal-Hitting-Time Scale}
\author{Yitzchak Shmalo}
\date{September 2026}

\begin{document}
\maketitle

\begin{abstract}
Ding, Lee and Peres asked whether, for finite connected graphs in the regime
$t_{\rm cov}/t_{\rm hit}\to\infty$, the standard deviation of the cover time is
universally of order at most the maximal hitting time, and whether the cover
time centered at $t_{\rm cov}$ has a universal constant-rate exponential tail
on the $t_{\rm hit}$ scale.  We give a simple family of unweighted graphs for
which both conclusions fail.  The graph $G_k$ consists of a center with $k$
pendant leaves and a $k$-clique attached through a single gateway edge.  We
prove the exact formula
\[
 t_{\rm hit}(G_k)=2k^2+k+1
\]
and $t_{\rm cov}(G_k)/t_{\rm hit}(G_k)=\Theta(\log k)$.  More sharply, for
\emph{every} sequence of starting vertices $x_k\in V(G_k)$,
\[
 \frac{\tau_{\rm cov}-t_{\rm cov}(G_k)}
 {t_{\rm hit}(G_k)\sqrt{(\log k)/2}}
 \Longrightarrow \N(0,1),
\]
and
\[
 \frac{\Var_{x_k}(\tau_{\rm cov})}
 {t_{\rm hit}(G_k)^2\log k}\longrightarrow \frac12.
\]
Thus $\operatorname{SD}(\tau_{\rm cov})/t_{\rm hit}\sim
\sqrt{(\log k)/2}\to\infty$.  The mechanism is a marked coupon collector:
while the leaves are collected, rare departures into the clique contribute
independent trap durations with variance of order $k^4$; their harmonic number
of occurrences produces the additional $\sqrt{\log k}$ fluctuation factor.
\end{abstract}

\section{The source question and the graph family}

Let $G$ be a finite connected graph and let simple random walk run on $G$.
Write $\tau_{\rm cov}$ for the first time every vertex has been visited,
\[
 t_{\rm cov}(G)=\max_{x\in V(G)}\E_x\tau_{\rm cov},
 \qquad
 t_{\rm hit}(G)=\max_{x,y\in V(G)}\E_xT_y.
\]
Ding--Lee--Peres \cite{DLP} ask in their Question 5.2, under the asymptotic
assumption $t_{\rm cov}(G_n)/t_{\rm hit}(G_n)\to\infty$, whether the standard
deviation of $\tau_{\rm cov}$ is bounded by $t_{\rm hit}$ and whether
$(\tau_{\rm cov}-t_{\rm cov})/t_{\rm hit}$ has exponential decay with a
constant rate.  We refute both statements on one explicit family.

For $k\ge3$, define $G_k$ as follows.  There is a center $s$, pendant leaves
$a_1,\dots,a_k$ adjacent to $s$, a clique
\[
 Q=\{g,q_2,\dots,q_k\}
\]
of size $k$, and one bridge edge $s g$.  Thus
\[
 |V(G_k)|=2k+1,
 \qquad
 |E(G_k)|=\frac{k^2+k+2}{2}.
\]
All edges are unweighted and the walk is the ordinary simple random walk.
Put
\[
 h_k=t_{\rm hit}(G_k).
\]

\begin{theorem}[Main theorem]\label{thm:main}
For $G_k$ as above,
\[
 h_k=2k^2+k+1,
 \qquad
 \frac{t_{\rm cov}(G_k)}{h_k}=\Theta(\log k).
\]
Moreover, for every choice of starting vertices $x_k\in V(G_k)$,
\begin{equation}\label{eq:clt-main}
 \frac{\tau_{\rm cov}-t_{\rm cov}(G_k)}
 {h_k\sqrt{(\log k)/2}}
 \Longrightarrow \N(0,1),
\end{equation}
and
\begin{equation}\label{eq:var-main}
 \frac{\Var_{x_k}(\tau_{\rm cov})}{h_k^2\log k}
 \longrightarrow \frac12.
\end{equation}
\end{theorem}

The source-level consequences are immediate.

\begin{corollary}\label{cor:refute}
Along $G_k$,
\[
 \frac{\operatorname{SD}_{x_k}(\tau_{\rm cov})}{t_{\rm hit}(G_k)}
 \sim \sqrt{\frac{\log k}{2}}\longrightarrow\infty
\]
uniformly over choices of starting vertices.  Hence there is no universal
constant $C$ such that
$\operatorname{SD}(\tau_{\rm cov})\le C t_{\rm hit}$ in the regime of
Question 5.2.

Furthermore, there do not exist universal constants $A,a>0$ for which
\[
 \Pp_x\!\left(\frac{\tau_{\rm cov}-t_{\rm cov}}{t_{\rm hit}}\ge u\right)
 \le A e^{-au}
 \qquad (u\ge0)
\]
for all such graphs and starts.  The same is therefore true for any stronger
two-sided version.
\end{corollary}

\begin{proof}
The standard-deviation assertion follows from \eqref{eq:var-main}.  For the
tail assertion fix, for example, $u_k=\frac14\sqrt{\log k}$.  By
\eqref{eq:clt-main},
\[
 \Pp_{x_k}\!\left(
 \frac{\tau_{\rm cov}-t_{\rm cov}}{h_k}\ge u_k
 \right)
 \longrightarrow 1-\Phi(1/(2\sqrt2))>0,
\]
whereas $Ae^{-a u_k}\to0$ for every fixed $A,a>0$.
\end{proof}

\section{Exact maximal hitting time}

We first calibrate the denominator exactly.  By symmetry there are only four
vertex types: the center $s$, a leaf, the gateway $g$, and a non-gateway clique
vertex $q$.

\begin{proposition}\label{prop:hit}
For every $k\ge3$,
\[
 t_{\rm hit}(G_k)=2k^2+k+1.
\]
\end{proposition}

\begin{proof}
We record the first-step solutions.  If $a$ is a fixed target leaf, $b\ne a$
is another leaf, and $q\in Q\setminus\{g\}$, then
\begin{align*}
 H(s,a)&=k^2+k+1,\\
 H(g,a)&=2k^2+2,\\
 H(q,a)&=2k^2+k+1,\\
 H(b,a)&=k^2+k+2.
\end{align*}
Indeed, writing $A=H(s,a)$, $B=H(g,a)$ and $C=H(q,a)$, the first-step
relations are
\[
 C=(k-1)+B,
 \qquad
 B=1+\frac{A}{k}+\frac{k-1}{k}C,
\]
and
\[
 A=1+\frac{k-1}{k+1}(1+A)+\frac{1}{k+1}B.
\]
Solving gives the displayed values.

For the other target types the corresponding values are
\begin{center}
\begin{tabular}{@{}lll@{}}
\toprule
Target & Start & Hitting time \\
\midrule
$s$ & $g$ & $k^2-k+1$ \\
$s$ & $q$ & $k^2$ \\
$s$ & leaf & $1$ \\
\addlinespace
$g$ & $s$ & $2k+1$ \\
$g$ & leaf & $2k+2$ \\
$g$ & $q$ & $k-1$ \\
\addlinespace
$q$ & $g$ & $k+3+4/k$ \\
$q$ & $s$ & $3k+4+4/k$ \\
$q$ & leaf & $3k+5+4/k$ \\
$q$ & $q'\ne q,g$ & $k+1+2/k$ \\
\bottomrule
\end{tabular}
\end{center}
These follow from the same two- or three-variable first-step systems.  For
example, with $X=H(g,q)$, $Y=H(q',q)$ and $Z=H(s,q)$,
\[
 Y=\frac{k-1+X}{2},\qquad
 Z=2k+1+X,\qquad
 X=1+\frac{Z}{k}+\frac{k-2}{k}Y,
\]
which yields $X=k+3+4/k$.  Comparing the table for $k\ge3$ shows that the
largest value is $H(q,a)=2k^2+k+1$.

As an independent consistency check, the commute-time identity counts the
same pair another way.  The effective resistance from a non-gateway clique
vertex $q$ to a leaf $a$ is
\[
 R_{\rm eff}(q,a)=\frac{2}{k}+1+1=2+\frac2k.
\]
Since $2|E|=k^2+k+2$,
\[
 H(q,a)+H(a,q)
 =2|E|R_{\rm eff}(q,a)
 =2k^2+4k+6+\frac4k,
\]
which agrees exactly with
$H(q,a)=2k^2+k+1$ and $H(a,q)=3k+5+4/k$.
\end{proof}

\section{The trap excursion}

A \emph{trap excursion} begins with a step $s\to g$ and ends at the next
return to $s$.  Let $Y_k$ be its duration, counting both the initial and final
bridge steps.

At $g$, the walk exits to $s$ with probability $1/k$.  On a failed exit it
moves to a non-gateway clique vertex and must return to $g$.  Let
\[
 F\sim\operatorname{Geom}_0(1/k)
\]
be the number of failed exits before success.  From a non-gateway clique
vertex, the hitting time $G$ of $g$ satisfies
\[
 G\sim\operatorname{Geom}_1(1/(k-1)).
\]
With iid $D_i=1+G_i$, independent of $F$,
\begin{equation}\label{eq:Yrep}
 Y_k=2+\sum_{i=1}^{F}D_i.
\end{equation}
Thus
\begin{equation}\label{eq:Ymom}
 \mu_k:=\E Y_k=k^2-k+2
\end{equation}
and, by the random-sum variance formula,
\begin{align}
 \sigma_k^2:=\Var(Y_k)
 &= (k-1)^2(k-2)+k^3(k-1)\notag\\
 &=k^4-4k^2+5k-2.\label{eq:Yvar}
\end{align}
We will also need a crude third moment.

\begin{lemma}\label{lem:Ythird}
There is a universal $C$ such that
\[
 \E Y_k^3\le Ck^6.
\]
\end{lemma}

\begin{proof}
The geometric variables in \eqref{eq:Yrep} satisfy
$\E F^3=O(k^3)$ and $\E D_1^3=O(k^3)$.  For nonnegative $d_i$,
\[
 \left(\sum_{i=1}^m d_i\right)^3
 \le m^2\sum_{i=1}^m d_i^3.
\]
Conditioning on $F$ therefore gives
\[
 \E\left(\sum_{i=1}^F D_i\right)^3
 \le \E[F^3]\E[D_1^3]=O(k^6).
\]
The deterministic additive constant $2$ only changes the universal constant.
\end{proof}

\section{Marked coupon collection}

Suppose the walk is at $s$ and exactly $r\in\{k-1,k\}$ specified leaves
remain to be visited.  Let $\Lambda_{k,r}$ be the additional time until all
those target leaves are visited.  When $j$ targets remain, discard from the
departure sequence all departures to already visited or non-target leaves.
Among the relevant symbols, the next symbol is
\begin{itemize}
\item a new target with probability $j/(j+1)$;
\item the trap edge $s g$ with probability $1/(j+1)$.
\end{itemize}
Consequently the number $N_j$ of trap departures before the next new target is
found satisfies
\[
 N_j\sim\operatorname{Geom}_0\!\left(\frac{j}{j+1}\right),
\]
and the variables $N_1,\dots,N_r$ are independent.  Their moments are
\begin{equation}\label{eq:Nmom}
 \E N_j=\frac1j,
 \qquad
 \Var(N_j)=\frac1j+\frac1{j^2},
 \qquad
 \E N_j^3=\frac1j+\frac6{j^2}+\frac6{j^3}\le\frac{13}{j}.
\end{equation}
Attach iid copies $Y_{k,j,\ell}$ of $Y_k$ to the trap symbols and set
\[
 Z_{k,j}=\sum_{\ell=1}^{N_j}Y_{k,j,\ell},
 \qquad
 S_{k,r}=\sum_{j=1}^r Z_{k,j}.
\]
The $Z_{k,j}$ are independent.  By \eqref{eq:Ymom}--\eqref{eq:Nmom},
\begin{align}
 b_{k,r}^2:=\Var(S_{k,r})
 &=(\sigma_k^2+\mu_k^2)H_r+\mu_k^2H_r^{(2)},\label{eq:bexact}
\end{align}
where $H_r=\sum_{j\le r}j^{-1}$ and
$H_r^{(2)}=\sum_{j\le r}j^{-2}$.  Hence for $r=k$ or $k-1$,
\begin{equation}\label{eq:basymp}
 b_{k,r}^2\sim 2k^4\log k.
\end{equation}

\begin{lemma}[Marked coupon-collector CLT]\label{lem:markedclt}
For either $r=k$ or $r=k-1$,
\[
 \frac{S_{k,r}-\E S_{k,r}}{b_{k,r}}
 \Longrightarrow \N(0,1).
\]
\end{lemma}

\begin{proof}
By nonnegativity,
\[
 Z_{k,j}^3
 \le N_j^2\sum_{\ell=1}^{N_j}Y_{k,j,\ell}^3.
\]
Conditioning on $N_j$, using Lemma \ref{lem:Ythird} and \eqref{eq:Nmom},
\[
 \E Z_{k,j}^3\le C\frac{k^6}{j}.
\]
For any nonnegative $X$,
$\E|X-\E X|^3\le8\E X^3$, so
\[
 \sum_{j=1}^r \E|Z_{k,j}-\E Z_{k,j}|^3
 =O(k^6\log k).
\]
By \eqref{eq:basymp}, $b_{k,r}^3=\Theta(k^6(\log k)^{3/2})$.
The Lyapunov ratio is therefore $O((\log k)^{-1/2})\to0$, proving the claim.
\end{proof}

We now put back the ordinary leaf excursions.  Delete the trap symbols from
the departure sequence and let $L_{k,r}$ be the number of leaf departures
needed to see all $r$ target leaves.  When $j$ targets remain the success
probability among leaf departures is $j/k$, so
\begin{equation}\label{eq:Lmom}
 \E L_{k,r}=kH_r,
 \qquad
 \Var(L_{k,r})
 =\sum_{j=1}^r\frac{1-j/k}{(j/k)^2}
 \le \frac{\pi^2}{6}k^2.
\end{equation}
All non-final leaf excursions have length two, while the final target is
covered on the first step into it.  Therefore, pathwise,
\begin{equation}\label{eq:lambda}
 \Lambda_{k,r}=S_{k,r}+2L_{k,r}-1.
\end{equation}
Since $k=o(b_{k,r})$, Lemma \ref{lem:markedclt}, \eqref{eq:Lmom} and Slutsky's
theorem imply
\begin{equation}\label{eq:lambdaclt}
 \frac{\Lambda_{k,r}-\E\Lambda_{k,r}}{b_{k,r}}
 \Longrightarrow\N(0,1),
\end{equation}
and, in fact,
\begin{equation}\label{eq:lambdavar}
 \frac{\Var(\Lambda_{k,r})}{b_{k,r}^2}\longrightarrow1.
\end{equation}
The mean is exact:
\begin{equation}\label{eq:lambdamean}
 \E\Lambda_{k,r}=(k^2+k+2)H_r-1.
\end{equation}
The variance conclusion follows because
$2(L_{k,r}-\E L_{k,r})/b_{k,r}\to0$ in $L^2$, not merely in probability.

\section{The source regime}

Starting from $s$, the graph cannot be covered before all $k$ leaves are
visited.  Hence \eqref{eq:lambdamean} gives
\[
 t_{\rm cov}(G_k)
 \ge (k^2+k+2)H_k-1.
\]
Matthews' bound gives, with $|V|=2k+1$,
\[
 t_{\rm cov}(G_k)\le h_k H_{2k}.
\]
Together with Proposition \ref{prop:hit},
\begin{equation}\label{eq:regime}
 \frac{t_{\rm cov}(G_k)}{h_k}=\Theta(\log k)\to\infty.
\end{equation}
Thus $G_k$ lies in precisely the asymptotic regime preceding DLP Question 5.2.

\section{From leaf completion to the actual cover time}

We prove that everything outside the marked coupon-collector clock is negligible
in $L^2$ on the scale $b_{k,r}$.

Fix an arbitrary start $x=x_k$.  Let $T_s$ be the first hitting time of $s$.
If $x$ is a leaf, put $r_k=k-1$; otherwise put $r_k=k$.  Before the first hit
of $s$, no pendant leaf other than the starting leaf can have been visited.
By the strong Markov property at $T_s$, define
\[
 V_k=T_s+\Lambda_{k,r_k}.
\]
Pathwise $\tau_{\rm cov}\ge V_k$.

\begin{lemma}\label{lem:Ts}
Uniformly in the starting vertex,
\[
 \E_x T_s^2=O(h_k^2)=O(k^4),
\]
and consequently
\[
 \frac{T_s-\E_xT_s}{b_{k,r_k}}\longrightarrow0
 \qquad\text{in }L^2.
\]
\end{lemma}

\begin{proof}
From every state $z$, $\E_zT_s\le h_k$.  Markov's inequality gives
$\Pp_z(T_s>2h_k)\le1/2$.  Applying the strong Markov property at deterministic
blocks of length $2h_k$ yields
\[
 \Pp_z(T_s>2mh_k)\le2^{-m}.
\]
Integration of this geometric tail gives, for example,
$\E_zT_s^2\le24h_k^2$.  Now use \eqref{eq:basymp}.
\end{proof}

Combining Lemma \ref{lem:Ts} with \eqref{eq:lambda} and
\eqref{eq:lambdavar},
\begin{equation}\label{eq:Vclt}
 \frac{V_k-\E V_k}{b_{k,r_k}}\Longrightarrow\N(0,1),
 \qquad
 \frac{\Var(V_k)}{b_{k,r_k}^2}\longrightarrow1.
\end{equation}

It remains to show that, by the time the leaves are complete, the clique is
also complete with probability $1-o(1)$, strongly enough for an $L^2$
transfer.

\begin{lemma}[One trap usually covers the clique]\label{lem:trapcovers}
Let $\delta_k$ be the probability that a trap excursion returns to $s$ before
visiting every vertex of $Q$.  Then
\[
 \delta_k=O\!\left(\frac{\log k}{k}\right).
\]
\end{lemma}

\begin{proof}
Couple the internal part of the excursion to simple random walk on $K_k$
started at $g$.  Set
$t_0=\lceil4(k-1)\log k\rceil$.  For any fixed clique vertex not equal to the
initial vertex, before it is hit each step hits it with probability
$1/(k-1)$.  A union bound therefore gives
\[
 \Pp(K_k\text{ is not covered by }t_0)
 \le k\left(1-\frac1{k-1}\right)^{t_0}
 \le k^{-3}.
\]
A bridge exit can occur only at a visit to $g$.  For the $K_k$ walk started at
$g$, for every $t\ge1$,
$\Pp(X_t=g)\le1/(k-1)$ (this follows directly from the two eigenvalues
$1$ and $-1/(k-1)$ of the complete-graph transition matrix).  Hence the
expected number of exit opportunities through time $t_0$ is at most
$2+4\log k$.  At each such opportunity the bridge is selected with probability
$1/k$.  Therefore
\[
 \delta_k\le k^{-3}+\frac{2+4\log k}{k}.
\]
\end{proof}

Let $B_k$ be the event that the clique is not fully covered at time $V_k$.
After $T_s$, restrict the departure sequence from $s$ to the trap symbol and
the $r_k$ still-unvisited target leaves.  These $r_k+1$ symbols are symmetric.
The probability that no trap occurs before every target leaf has appeared is
exactly the probability that the trap is last in first-appearance order,
namely $1/(r_k+1)$.  If a trap does occur and its first excursion covers the
clique, then $B_k$ is impossible.  Lemma \ref{lem:trapcovers} therefore gives
\begin{equation}\label{eq:Bprob}
 \Pp_x(B_k)
 \le \frac1{r_k+1}+\delta_k
 =O\!\left(\frac{\log k}{k}\right)
\end{equation}
uniformly over $x$.

Set
\[
 D_k=\tau_{\rm cov}-V_k\ge0.
\]
Then $D_k=0$ on $B_k^c$.  We next need a uniform second-moment bound for a
fresh full cover time.  From every state $z$,
$\E_z\tau_{\rm cov}\le t_{\rm cov}$.  The same block-Markov argument as in
Lemma \ref{lem:Ts} yields
\begin{equation}\label{eq:coversecond}
 \sup_z\E_z\tau_{\rm cov}^2\le24t_{\rm cov}^2.
\end{equation}
At the stopping time $V_k$, the additional time to cover the still-missing
vertices is stochastically bounded by the time to cover the whole graph afresh
from the current state.  Hence, using \eqref{eq:Bprob},
\begin{align}
 \E_xD_k^2
 &\le 24\Pp_x(B_k)t_{\rm cov}(G_k)^2\notag\\
 &=O(k^3(\log k)^3),\label{eq:D2}
\end{align}
because Matthews' bound and Proposition \ref{prop:hit} give
$t_{\rm cov}=O(k^2\log k)$.  Comparing \eqref{eq:D2} with
$b_{k,r_k}^2\sim2k^4\log k$ shows
\begin{equation}\label{eq:Dsmall}
 \frac{D_k-\E D_k}{b_{k,r_k}}\longrightarrow0
 \qquad\text{in }L^2.
\end{equation}

\section{Worst-start centering and proof of the main theorem}

For any two starts $x,y$,
\[
 \E_y\tau_{\rm cov}\le H(y,x)+\E_x\tau_{\rm cov}\le h_k+\E_x\tau_{\rm cov}.
\]
Taking $y$ to attain the maximum in the definition of $t_{\rm cov}$ yields
\begin{equation}\label{eq:meanstart}
 0\le t_{\rm cov}(G_k)-\E_x\tau_{\rm cov}\le h_k.
\end{equation}
Since $h_k=o(b_{k,r_k})$, the source's worst-start centering differs
negligibly from the mean under any fixed start at our fluctuation scale.

Now
\[
 \tau_{\rm cov}-\E_x\tau_{\rm cov}
 =(V_k-\E V_k)+(D_k-\E D_k).
\]
Equations \eqref{eq:Vclt} and \eqref{eq:Dsmall} imply
\begin{equation}\label{eq:tau-b-clt}
 \frac{\tau_{\rm cov}-\E_x\tau_{\rm cov}}{b_{k,r_k}}
 \Longrightarrow\N(0,1),
 \qquad
 \frac{\Var_x(\tau_{\rm cov})}{b_{k,r_k}^2}\longrightarrow1,
\end{equation}
uniformly over the four possible starting-orbit types.  By
\eqref{eq:meanstart}, the first convergence remains true with
$t_{\rm cov}(G_k)$ in place of $\E_x\tau_{\rm cov}$.

Finally, \eqref{eq:basymp} and Proposition \ref{prop:hit} give
\[
 \frac{b_{k,r_k}^2}{h_k^2\log k}
 \longrightarrow
 \frac{2k^4}{(2k^2)^2}
 =\frac12.
\]
Substituting this into \eqref{eq:tau-b-clt} proves
\eqref{eq:clt-main} and \eqref{eq:var-main}, completing the proof of
Theorem \ref{thm:main}.

\section{Interpretation}

The obstruction is not an unusually large expected cover time.  In fact
$t_{\rm cov}/t_{\rm hit}$ grows only logarithmically, exactly as for many
coupon-collector-type examples.  The obstruction is that each coupon-collector
stage has a small chance of triggering a long clique excursion whose standard
deviation is already of order $k^2$.  Summing those independent marked stages
over the harmonic sequence $j=1,\dots,k$ creates variance
$\asymp k^4\log k$, while the maximal hitting time is only $\asymp k^2$.

This also explains why results under vertex-transitivity or strong spatial
decorrelation do not extend automatically to arbitrary graphs.  The graph
$G_k$ deliberately couples a coupon collector to a metastable region through a
single gateway.  The cover state may be almost complete while the durations of
covered-region detours continue to carry macroscopic random variance.

\section*{Prior-art status}

The original source is \cite{DLP}.  Ding proved geometric-mean-scale cover-time
concentration for general trees and bounded-degree settings \cite{Ding}; Zhai
proved a general GFF-based exponential concentration theorem \cite{Zhai};
Hermon characterized qualitative concentration for vertex-transitive graphs
and a wider reversible class \cite{Hermon}; and Berestycki--Hermon--Teyssier
obtained Gumbel fluctuation results under bounded-degree vertex-transitive
hypotheses \cite{BHT}, with refined hitting-time estimates in a 2026 companion
preprint \cite{BHT26}.  A bounded literature search dated 19 September 2026
did not locate the unrestricted counterexample above.  This negative search is
\emph{not} a proof of novelty or priority.  The detailed search record is kept
beside this manuscript in `prior_art_2026-09-19.md`.

\begin{thebibliography}{9}
\bibitem{DLP}
J. Ding, J. R. Lee, and Y. Peres,
\emph{Cover times, blanket times, and majorizing measures},
Ann. of Math. 175 (2012), 1409--1471.
\url{https://doi.org/10.4007/annals.2012.175.3.8}

\bibitem{Ding}
J. Ding,
\emph{Asymptotics of cover times via Gaussian free fields: bounded-degree graphs and general trees},
arXiv:1103.4402.

\bibitem{Zhai}
A. Zhai,
\emph{Exponential concentration of cover times},
arXiv:1407.7617; Electron. J. Probab. (2018).

\bibitem{Hermon}
J. Hermon,
\emph{A spectral characterization for concentration of the cover time},
J. Theoret. Probab. 33 (2020), 2167--2184.

\bibitem{BHT}
N. Berestycki, J. Hermon, and L. Teyssier,
\emph{On the universality of fluctuations for the cover time},
arXiv:2202.02255.

\bibitem{BHT26}
N. Berestycki, J. Hermon, and L. Teyssier,
\emph{Stationary hitting times on vertex-transitive graphs},
arXiv:2601.03864 (2026).
\end{thebibliography}

\end{document}
